\section{Datenbanken und Tabellen anlegen und bearbeiten (Überblick) \tagLK}

Bisher haben wir Daten aus einer bereits fertigen Datenbank ausgelesen
(\lstinline|SELECT|). Im Folgenden erhältst du einen kompakten Überblick über
die wichtigsten SQL-Befehle zum Anlegen und Verändern von Datenbankstrukturen
sowie zum Einfügen, Ändern und Löschen von Datensätzen.

\noindent \textbf{Hinweis für den Leistungskurs:} Auch wenn wir hier zunächst
keine eigenen Datenbanken von Grund auf neu entwerfen, ist ein sicheres
Verständnis von \lstinline|CREATE TABLE| essenziell, um bestehende
Datenbankskripte (wie \myFile{haustier.sql}) lesen und nachvollziehen zu
können. Darüber hinaus wirst du die Befehle \lstinline|INSERT INTO|,
\lstinline|UPDATE| und \lstinline|DELETE| in Zukunft noch regelmäßig benötigten
-- insbesondere beim schreibenden Zugriff auf Datenbanken aus Java heraus.

\begin{framed}
\textbf{Merksatz (DDL und DML):} SQL-Befehle lassen sich nach ihrer Aufgabe in
zwei Kategorien unterteilen:
\begin{compactitem}
  \item \textbf{DDL (\emph{Data Definition Language} -- Datendefinitionssprache):}
  Befehle, die die Struktur (den „Bauplan“) einer Datenbank festlegen, ändern
  oder löschen (\lstinline|CREATE|, \lstinline|ALTER|, \lstinline|DROP|).
  \item \textbf{DML (\emph{Data Manipulation Language} -- Datenmanipulationssprache):}
  Befehle, die die gespeicherten Daten (den Inhalt der Tabellen) abfragen oder
  verändern (\lstinline|SELECT|, \lstinline|INSERT|, \lstinline|UPDATE|,
  \lstinline|DELETE|).
\end{compactitem}
\end{framed}

\subsection{Datenbanken anlegen und auswählen}

Eine neue Datenbank (ein neues Schema) legt man mit \lstinline|CREATE SCHEMA|
an. Dabei sollte man direkt den Zeichensatz festlegen (wir verwenden immer
UTF-8, in MySQL modern als \lstinline|utf8mb4| bezeichnet):

\begin{lstlisting}
CREATE SCHEMA `\slshape{datenbankname}` DEFAULT CHARACTER SET utf8mb4;
\end{lstlisting}

Mit \lstinline|SHOW SCHEMAS;| lassen sich alle vorhandenen Datenbanken
anzeigen. Damit sich alle nachfolgenden Anweisungen auf die gewünschte
Datenbank beziehen, aktiviert man sie mit:

\begin{lstlisting}
USE `\slshape{datenbankname}`;
\end{lstlisting}


\subsection{Tabellen erstellen}

\begin{lstlisting}
CREATE TABLE [IF NOT EXISTS] `\slshape{tabellenname}`
(
  `\slshape{Spalten-Definition1}`,
  ...  
  `\slshape{Spalten-DefinitionN}`
);
\end{lstlisting}

Ausdrücke in eckigen Klammern sind optional. Bei Angabe von
\lstinline|IF NOT EXISTS| wird die Tabelle nur dann angelegt, wenn sie noch
nicht existiert. Eine \textit{Spalten-Definition} hat folgendes Schema:

\begin{lstlisting}
`\slshape{spaltenname} typ` [NOT NULL] [DEFAULT `\slshape{wert}`] [AUTO_INCREMENT]
\end{lstlisting}

\begin{compactitem}
  \item \textit{Typ} legt den Datentyp der Spalte fest (siehe
  Tabelle~\ref{tab:mysql-datentypen}).
  \item \lstinline|NOT NULL| legt fest, dass das Feld nicht leer
  (\lstinline|NULL|) bleiben darf.
  \item \lstinline|DEFAULT| bestimmt einen Standardwert für neue Datensätze.
  \item \lstinline|AUTO_INCREMENT| nummeriert eine Ganzzahlspalte automatisch
  durch (beginnend bei $1$). Dies ist nur für maximal eine indizierte Spalte
  pro Tabelle erlaubt -- typischerweise den Primärschlüssel \linebreak
  (\lstinline|PRIMARY KEY (|\textit{spaltenname}\lstinline|)|).
\end{compactitem}

\begin{table}[h]
\centering
\begin{tabular}{|l|p{11cm}|}
\hline
{\bfseries Typ} & {\bfseries Bedeutung}\\
\hline
\texttt{INT} & 4-Byte großer Ganzzahlwert\\
\hline
\texttt{FLOAT} & Fließkommazahl mit einfacher Genauigkeit\\
\hline
\texttt{DOUBLE} & Fließkommazahl mit doppelter Genauigkeit\\
\hline
\texttt{VARCHAR(\slshape{maximale Länge})} &
Zeichenkette. In Klammern gibt man die maximale Anzahl der Zeichen an (in
neueren MySQL-Versionen bis zu $65.535$ Zeichen).\\
\hline
\texttt{ENUM(\slshape{Wert1, Wert2, \ldots})}
& ermöglicht die Aufzählung einer Werteliste, z.\,B.:\\
& \texttt{ENUM('rot', 'blau', 'grün')}\\
\hline
\texttt{DATE} & Datum im Format YYYY-MM-DD\\
\hline
\texttt{TIME} & Uhrzeit im Format HH:MM:SS\\
\hline
\texttt{DATETIME} & Datum und Uhrzeit im Format YYYY-MM-DD HH:MM:SS\\
\hline
\end{tabular}
\caption{Datentypen von MySQL (Auswahl)}
\label{tab:mysql-datentypen}
\end{table}

Den Aufbau einer bestehenden Tabelle kann man jederzeit mit
\lstinline|DESCRIBE| \textit{tabellenname}\lstinline|;| überprüfen; alle Tabellen einer
Datenbank zeigt \lstinline|SHOW TABLES;| an.


\subsection{Datenbanken und Tabellen löschen}

\begin{lstlisting}
DROP SCHEMA [IF EXISTS] `\slshape{datenbankname}`;
DROP TABLE [IF EXISTS] `\slshape{tabellenname}`;
\end{lstlisting}


\subsection{Tabellenstruktur ändern}

Mit der Anweisung \lstinline|ALTER TABLE| \ldots\ kann man die Struktur einer
existierenden Tabelle nachträglich verändern -- etwa Spalten hinzufügen oder
entfernen:

\begin{lstlisting}
ALTER TABLE `\slshape{tabellenname}` ADD COLUMN (`\slshape{Spalten-Definition1, Spalten-Definition2, \ldots}`); 
ALTER TABLE `\slshape{tabellenname}` DROP COLUMN `\slshape{spaltenname}`;
\end{lstlisting}

Es können auch mehrere Änderungen in einer \lstinline|ALTER|-Anweisung
zusammengefasst werden -- wie in diesem Beispiel:

\begin{lstlisting}[numbers=left, xleftmargin=7mm]
DROP SCHEMA IF EXISTS stammbaum;
CREATE SCHEMA stammbaum DEFAULT CHARACTER SET utf8mb4;
USE stammbaum;
CREATE TABLE person (
  vorname VARCHAR(20) NOT NULL,
  nachname VARCHAR(20) NOT NULL,
  geschlecht ENUM('männlich', 'weiblich') DEFAULT 'männlich', 
  anzahl_kinder INT NOT NULL DEFAULT 0
);
INSERT INTO person VALUES
('Hartmut', 'Meyer', 'männlich', 3);
ALTER TABLE person
  DROP COLUMN anzahl_kinder,
  ADD COLUMN (
    person_id INT AUTO_INCREMENT,
    anzahl_toechter INT NOT NULL DEFAULT 0,
    anzahl_soehne INT NOT NULL DEFAULT 0
  ),
  ADD PRIMARY KEY (person_id)
;
UPDATE person SET anzahl_toechter=1, anzahl_soehne=2
WHERE vorname='Hartmut' AND nachname='Meyer';
\end{lstlisting}




